
\section{上同调}\label{sec:D-6}

\BookRecordIndex{T0023}Grothendieck--Knudsen 流形 $\overline{\mathcal M}_{0,n}$ 有许多值得注意的性质。定理~\ref{thm:D-4-7} 已表明，每个 \BookRecordIndex{T0008}$n$-分割网络都确定一个复子流形。这些子流形生成 $\overline{\mathcal M}_{0,n}$ 的上同调，并提供关于其环结构的有用信息。在下面的讨论中，有时将 $\overline{\mathcal M}_{0,n}$ 看作装饰树的集合，有时则将其视为带坐标 $w\in(S^2)^N$ 的流形 $\overline M_{0,n}$，会比较方便。

\paragraph{普适曲线。}
存在自然投影
\begin{equation}\label{eq:D-6-1}
  \pi:\overline M_{0,n+1}\longrightarrow\overline M_{0,n},
\end{equation}
它只保留不含指标 $n+1$ 的坐标 $w_{ijk\ell}$。该映射显然全纯，且取值于 $\overline M_{0,n}$（见式~\eqref{eq:D-4-4}）。在定理~\ref{thm:D-4-2} 的双射 $\overline{\mathcal M}_{0,n+1}\cong\overline M_{0,n+1}$ 下，它对应于删除第 $n+1$ 个\BookRecordIndex{T0066}标记点。在几何上可以这样描述：若承载第 $n+1$ 个\BookRecordIndex{T0066}标记点的顶点 $\alpha$ 有至少四个\BookRecordIndex{T0320}特殊点，删除该\BookRecordIndex{T0066}标记点后所得曲线仍稳定。否则，该顶点恰有三个\BookRecordIndex{T0320}特殊点，其中至少一个是\BookRecordIndex{T0191}节点。\sourcefootnote{原书实体页 658 此处写作 ``exactly three marked points, at least one of them nodal''；\BookRecordIndex{T0191}节点并非\BookRecordIndex{T0066}标记点。由本句后续的两种\BookRecordIndex{T0191}节点情形及稳定性定义，将 ``marked points'' 改作“\BookRecordIndex{T0320}特殊点”。}若它有两个\BookRecordIndex{T0191}节点，就删除该顶点，改以连接 $\alpha$ 的两条边所连的两个顶点上的一对新\BookRecordIndex{T0191}节点代替。若它有一个\BookRecordIndex{T0191}节点，以及指标为 $m$ 和 $n+1$ 的两个\BookRecordIndex{T0066}标记点，就删除该顶点，改在与 $\alpha$ 以一条边相连的唯一顶点上安放第 $m$ 个\BookRecordIndex{T0066}标记点。

对 $i=1,\ldots,n$，令 $Z_i\subset\overline M_{0,n+1}$ 为对应于 $(n+1)$-分割 $I_i:=\{1,\ldots,n\}\setminus\{i\}$ 的子集。它由所有带 $n+1$ 个\BookRecordIndex{T0066}标记点的亏格零\BookRecordIndex{T0339}稳定曲线同构类组成，其中指标 $i$ 与 $n+1$ 的\BookRecordIndex{T0066}标记点同处一个顶点，而且该顶点恰有三个\BookRecordIndex{T0320}特殊点。定理~\ref{thm:D-4-2} 的双射将 $Z_i$ 映为
\begin{equation}\label{eq:D-6-2}
  Z_i:=\bigl\{w\in\overline M_{0,n+1}\mid
  w_{jj'in+1}=\infty\quad
  \forall j,j'\in\{1,\ldots,n\}\setminus\{i\}\bigr\}.
\end{equation}
由定理~\ref{thm:D-4-7}，这是 $\overline M_{0,n+1}$ 的复子流形，且投影~\eqref{eq:D-6-1} 在其上的限制是微分同胚 $\pi:Z_i\to\overline M_{0,n}$。其逆是 $\pi$ 的截面，记作
\begin{equation}\label{eq:D-6-3}
  \sigma_i:\overline M_{0,n}\longrightarrow\overline M_{0,n+1},
  \qquad \pi\circ\sigma_i=\mathrm{id},\qquad
  \sigma_i(\overline M_{0,n})=Z_i,\quad i=1,\ldots,n.
\end{equation}
显式地，$\sigma_i$ 将 $w\in\overline M_{0,n}$ 送到如下定义的 $\sigma_i(w)\in\overline M_{0,n+1}$：
\[
  (\sigma_i(w))_{jk\ell m}:=
  \begin{cases}
    w_{jk\ell m},&j,k,\ell,m\le n,\\
    w_{jk\ell i},&m=n+1,\quad j,k,\ell\ne i,\\
    w_{jkim},&\ell=n+1,\quad j,k,m\ne i,\\
    w_{ji\ell m},&k=n+1,\quad j,\ell,m\ne i,\\
    w_{ik\ell m},&j=n+1,\quad k,\ell,m\ne i,\\
    0,&\{j,m\}=\{i,n+1\}\text{ 或 }\{k,\ell\}=\{i,n+1\},\\
    1,&\{j,\ell\}=\{i,n+1\}\text{ 或 }\{k,m\}=\{i,n+1\},\\
    \infty,&\{j,k\}=\{i,n+1\}\text{ 或 }\{\ell,m\}=\{i,n+1\}.
  \end{cases}
\]
这里 $j,k,\ell,m\in\{1,\ldots,n+1\}$ 两两互异。从几何上看，截面 $\sigma_i$ 把第 $i$ 个\BookRecordIndex{T0066}标记点“加倍”：它在树中引入一个与 $\alpha_i$ 以边相连的新顶点，该顶点承载指标 $i$ 与 $n+1$ 的\BookRecordIndex{T0066}标记点。

\begin{exercise}\label{ex:D-6-1}
用建模在\BookRecordIndex{T0084}带标号树上的 $SC_{0,n}$ 元组，写出投影~\eqref{eq:D-6-1} 和截面~\eqref{eq:D-6-3} 的显式公式。
\end{exercise}

\begin{exercise}\label{ex:D-6-2}
证明映射 $\pi:\overline M_{0,n+1}\to\overline M_{0,n}$ 是一个\BookRecordIndex{T0191}节点族，即 $\pi$ 的每个\BookRecordIndex{T0233}临界点都是\BookRecordIndex{T0191}节点（见文献~\srcite{337} 的定义 4.2）。证明 $M_{0,n+1}$ 中各点均为 $\pi$ 的\BookRecordIndex{T0405}正则点，$M_{0,n}$ 是 $\pi$ 的\BookRecordIndex{T0409}正则值集（见式~\eqref{eq:D-4-11}）。\sourcefootnote{原书实体页 659 写作 ``$M_{0,n+1}$ is the set of regular points''。这不能表示\BookRecordIndex{T0405}正则点集恰等于 $M_{0,n+1}$：例如 $n=3$ 时底空间 $\overline M_{0,3}$ 是单点，$\pi$ 的全部点均正则，而 $M_{0,4}$ 不含边界点。此处只将原句改为可成立的包含断言。}
\end{exercise}

由习题~\ref{ex:D-6-2} 可知，$\pi:\overline M_{0,n+1}\to\overline M_{0,n}$ 的每条纤维都是带 $n$ 个\BookRecordIndex{T0066}标记点的\BookRecordIndex{T0191}节点黎曼曲面，其\BookRecordIndex{T0066}标记点由截面 $Z_1,\ldots,Z_n$ 给出。事实上，对每条\BookRecordIndex{T0339}稳定曲线 $\mathbf z\in SC_{0,n}$，在式~\eqref{eq:D-4-1} 定义的点 $w(\mathbf z)\in\overline M_{0,n}$ 上方的纤维，同构于式~\eqref{eq:D-3-2} 定义的空间 $\Sigma(\mathbf z)$ 本身：
\begin{equation}\label{eq:D-6-4}
  \pi^{-1}(w(\mathbf z))\cong\Sigma(\mathbf z),\qquad
  \pi^{-1}(w(\mathbf z))\cap Z_i\cong[\alpha_i,z_i].
\end{equation}
为说明这一点，若 $[\alpha,z]$ 既不是 $\Sigma(\mathbf z)$ 的\BookRecordIndex{T0191}节点，也不等于 $n$ 个\BookRecordIndex{T0066}标记点中任何一个，就可把它作为第 $n+1$ 个\BookRecordIndex{T0066}标记点加入，得到纤维中的一个元素。若它恰为\BookRecordIndex{T0191}节点或某个\BookRecordIndex{T0066}标记点，则可添加一个球面，得到纤维中的元素。

带 $n$ 个截面、满足如下条件的\BookRecordIndex{T0191}节点族 $(\pi_B:Q\to B,Z_1,\ldots,Z_n)$ 称为\emph{带 $n$ 个\BookRecordIndex{T0066}标记点的亏格零\BookRecordIndex{T0191}节点黎曼曲面普适族}（见文献~\srcite{337} 的定义 5.1 与 6.2）。

\noindent\textup{（I）}对任意另一\BookRecordIndex{T0191}节点族 $(\pi_A:P\to A,R_1,\ldots,R_n)$ 及任意纤维同构 $f:P_a\to Q_b$，若 $f(P_a\cap R_i)=Q_b\cap Z_i$ 对 $i=1,\ldots,n$ 成立，则 $f$ 在 $P$ 中 $P_a$ 的某个邻域上唯一地延拓为\BookRecordIndex{T0191}节点族态射。

\noindent\textup{（II）}$Q$ 的每条纤维都是带 $n$ 个\BookRecordIndex{T0066}标记点的亏格零\BookRecordIndex{T0339}稳定曲线；反之，每条带 $n$ 个\BookRecordIndex{T0066}标记点的亏格零\BookRecordIndex{T0339}稳定曲线都恰与其中一条纤维同构。

文献~\srcite{337} 的定理 5.4 表明，带 $n$ 个截面的\BookRecordIndex{T0191}节点族 $(\pi_B:Q\to B,Z_1,\ldots,Z_n)$ 满足（I），当且仅当与每条纤维关联的线性化\BookRecordIndex{T0210}柯西--黎曼算子为双射。该算子的定义域是沿纤维的向量场空间，其中向量场投影到基底的一个常值切向量，且与截面 $Z_i$ 相切；值域是纤维上取值于\BookRecordIndex{T0308}竖直切丛的 $(0,1)$-形式空间。由~\eqref{eq:D-6-1}、\eqref{eq:D-6-2} 定义的\BookRecordIndex{T0191}节点族 $(\pi:\overline M_{0,n+1}\to\overline M_{0,n},Z_1,\ldots,Z_n)$ 按定义满足（II）；它还满足（I），这正是下一个习题的内容。它通常称为 $\overline M_{0,n}$ 上的普适曲线。

\begin{exercise}\label{ex:D-6-3}
证明 $(\pi:\overline M_{0,n+1}\to\overline M_{0,n},Z_1,\ldots,Z_n)$ 是带 $n$ 个\BookRecordIndex{T0066}标记点的亏格零\BookRecordIndex{T0191}节点黎曼曲面普适族。
\end{exercise}

\paragraph{上同调的生成元。}
Keel 在文献~\srcite{208} 中证明，与 \BookRecordIndex{T0008}$n$-分割网络 $S\in\mathcal S_n$ 相应的复子流形 $V_S$ 在加法意义上生成 $\overline M_{0,n}$ 的上同调。

\begin{theorem}[Keel]\label{thm:D-6-4}
$\overline M_{0,n}$ 单连通，且 $H^{\mathrm{odd}}(\overline M_{0,n};\mathbb Z)=\{0\}$。此外，若 $a\in H^{2k}(\overline M_{0,n};\mathbb Z)$ 对每个满足 $\#S=n-3-k$ 的 \BookRecordIndex{T0008}$n$-分割网络 $S\in\mathcal S_n$ 都有 $\langle a,[V_S]\rangle=0$，则 $a=0$。
\end{theorem}

\begin{proof}[单连通性的证明]
$\overline M_{0,3}=\{\mathrm{pt}\}$ 显然单连通。令 $n\ge4$，归纳假设 $\overline M_{0,n-1}$ 单连通。取光滑圈 $\gamma:\mathbb R/\mathbb Z\to\overline M_{0,n}$。由标准\BookRecordIndex{T0166}横截性理论，可设 $\gamma$ 与每个 \BookRecordIndex{T0008}$n$-分割 $I\in\mathcal I_n$ 的 $V_I$ 横截，因而 $\gamma(t)$ 对所有 $t$ 均处于\BookRecordIndex{T0106}顶层 $M_{0,n}$。现在选圈的同伦
\[
  \mathbb R/\mathbb Z\times[0,1]\longrightarrow\overline M_{0,n},
  \qquad(t,\lambda)\longmapsto\gamma_\lambda(t),
\]
使 $i,j,k,\ell\le n-1$ 时 $w_{ijk\ell}(\gamma_\lambda(t))=w_{ijk\ell}(\gamma(t))$，且
\[
  w_{123n}(\gamma_\lambda(t))
  =(1-\lambda)w_{123n}(\gamma(t))
    +\lambda w_{123,n-1}(\gamma(t)).
\]
这些条件唯一确定 $\gamma_\lambda$，并保证 $\gamma_0(t)=\gamma(t)$，而 $\gamma_1(t)$ 位于投影 $\pi:\overline M_{0,n}\to\overline M_{0,n-1}$ 的一个光滑截面的像中。几何上，$\gamma_0$ 位于 $\overline M_{0,n}$ 的一个子集内，该子集可与 $M_{0,n-1}$ 上的一个复线丛等同（即 $i,j,k,\ell\le n-1$ 时 $w_{ijk\ell}\notin\{0,1,\infty\}$ 且 $w_{123n}\ne\infty$ 的子集），因此 $\gamma_0$ 同伦于这个线丛零截面中的圈。归纳假设说 $\overline M_{0,n-1}$ 单连通，故 $\gamma_1$ 可缩。这说明 $\overline M_{0,n}$ 单连通。
\end{proof}

Keel 定理的完整证明涉及吹起构造，这里不讨论该论证。若能找到\BookRecordIndex{T0246}模空间 $\overline M_{0,n}$ 上完美 Morse 函数的显式公式，将会很有意思。我们不知道已有文献中存在这样的构造。完美 Morse 函数的存在，对 $n=3,4,5$ 可由 D.7 节的讨论\BookRecordIndex{T0326}推出；对 $n\ge6$，则由证明 $h$-配边定理时使用的标准手柄消去过程\BookRecordIndex{T0326}推出。

\paragraph{\BookRecordIndex{T0378}引力后裔。}
对 $i=1,\ldots,n$，有复线丛 $L_i\to\overline M_{0,n}$，定义为在截面 $\sigma_i:\overline M_{0,n}\to Z_i$ 下拉回 $Z_i\subset\overline M_{0,n+1}$ 的余\BookRecordIndex{T0117}法丛。从几何上看，$L_i$ 在 $w(\mathbf z)$ 上的纤维可视为 $\Sigma(\mathbf z)$ 在第 $i$ 个\BookRecordIndex{T0066}标记点的余切空间。\sourcefootnote{原书实体页 660 在定义 $L_i$ 为截面 $Z_i$ 的余\BookRecordIndex{T0117}法丛拉回后，把纤维称为 ``tangent space''；余\BookRecordIndex{T0117}法丛的纤维应为纤维方向的余切空间。此处据紧邻定义补正。}记此线丛的第一陈类为
\begin{equation}\label{eq:D-6-5}
  \psi_i:=c_1(L_i)\in H^2(\overline M_{0,n};\mathbb Z).
\end{equation}
这些上同调类是所谓\emph{\BookRecordIndex{T0378}引力后裔}的最简单特例。以下结论由 Givental 在文献~\srcite{149} 中证明。

\begin{theorem}\label{thm:D-6-5}
设非负整数 $d_1,\ldots,d_n$ 满足 $d_1+\cdots+d_n=n-3$。则
\begin{equation}\label{eq:D-6-6}
  \int_{\overline M_{0,n}}\psi_1^{d_1}\smile\psi_2^{d_2}
    \smile\cdots\smile\psi_n^{d_n}
  =\frac{(n-3)!}{d_1!\cdots d_n!}.
\end{equation}
\end{theorem}
诸式~\eqref{eq:D-6-6} 可合写为一个等式
\begin{equation}\label{eq:D-6-7}
  \int_{\overline M_{0,n}}\prod_{i=1}^{n}\frac{1}{1-t_i\psi_i}
  =(t_1+\cdots+t_n)^{n-3}.
\end{equation}
等式两边都是变量 $t_1,\ldots,t_n$ 的形式幂级数，因此比较系数可知，\eqref{eq:D-6-6} 与~\eqref{eq:D-6-7} 等价。

定理~\ref{thm:D-6-5} 的证明基于 $\psi_i$ 类的递推公式。对 $n\ge3$ 和 $i=1,\ldots,n$，分别记投影~\eqref{eq:D-6-1} 为 $\pi_{n+1}:\overline M_{0,n+1}\to\overline M_{0,n}$，截面~\eqref{eq:D-6-3} 为 $\sigma_{i,n}:\overline M_{0,n}\to\overline M_{0,n+1}$，其像为 $Z_{i,n+1}\subset\overline M_{0,n+1}$，$\sigma_{i,n}$ 拉回 $Z_{i,n+1}$ 的余\BookRecordIndex{T0117}法丛为 $L_{i,n}\to\overline M_{0,n}$，其第一陈类~\eqref{eq:D-6-5} 为 $\psi_{i,n}\in H^2(\overline M_{0,n};\mathbb Z)$，并记
\[
  e_{i,n+1}:=\operatorname{PD}([Z_{i,n+1}])
     \in H^2(\overline M_{0,n+1};\mathbb Z)
\]
为子流形 $Z_{i,n+1}$ 的庞加莱对偶。于是 $\psi_{i,n}=-\sigma_{i,n}^{*}e_{i,n+1}$。

\begin{proof}[定理~\ref{thm:D-6-5} 的证明]
证明取自文献~\srcite{149}，分三步。

\emph{第一步。} 对 $i=1,\ldots,n$ 及 $n\ge3$，上同调类 $\psi_{i,n+1}$ 与 $\psi_{i,n}$ 满足递推公式
\begin{equation}\label{eq:D-6-8}
  \psi_{i,n+1}=\pi_{n+1}^{*}\psi_{i,n}+e_{i,n+1}.
\end{equation}
下面简述证明。线丛 $L_{i,n+1}$ 与 $\pi_{n+1}^{*}L_{i,n}$ 在 $\overline M_{0,n+1}\setminus Z_{i,n+1}$ 上同构。因此它们的第一陈类相差 $e_{i,n+1}$ 的某个整数倍；也就是说，存在整数 $k$，使 $\psi_{i,n+1}-\pi_{n+1}^{*}\psi_{i,n}=k e_{i,n+1}$。我们须证 $k=1$。注意 $L_{i,n+1}$ 限制到 $Z_{i,n+1}$ 后可平凡化，而 $\pi_{n+1}^{*}L_{i,n}$ 限制到 $Z_{i,n+1}$ 后是余\BookRecordIndex{T0117}法丛。因此 $\psi_{i,n+1}-\pi_{n+1}^{*}\psi_{i,n}$ 在 $Z_{i,n+1}$ 上的限制是\BookRecordIndex{T0117}法丛的欧拉类，从而等于 $e_{i,n+1}$ 在 $Z_{i,n+1}$ 上的限制。该类非零，故 $k=1$，第一步得证。\sourcefootnote{原书此步称限制类非零，但在所述范围的端点 $n=3$，$Z_{i,4}\cong\overline M_{0,3}$ 是单点，其 $H^2$ 为零；因此这一限制论证不能单独确定该端点的 $k$。此处保留原论证，提示端点须另行核验。}

\emph{第二步。} 若 $d_i\ge0$ 且 $d_1+\cdots+d_n=n-2$，则
\begin{equation}\label{eq:D-6-9}
  \int_{\overline M_{0,n+1}}\prod_{i=1}^{n}\psi_{i,n+1}^{d_i}
  =\sum_{d_i>0}\int_{\overline M_{0,n}}
    \psi_{i,n}^{d_i-1}\prod_{j\ne i}\psi_{j,n}^{d_j}.
\end{equation}
由于线丛 $L_{i,n+1}\to\overline M_{0,n+1}$ 在 $Z_{i,n+1}$ 上可平凡化，其第一陈类 $\psi_{i,n+1}$ 是一个与 $Z_{i,n+1}$ 不相交的子流形的庞加莱对偶。因此 $\psi_{i,n+1}\smile e_{i,n+1}=0$，再由~\eqref{eq:D-6-8} 得
\[
  \pi_{n+1}^{*}\psi_{i,n}^{d}
  =\psi_{i,n+1}^{d}+(-e_{i,n+1})^{d},\qquad d>0.
\]
于是
\begin{align*}
  0
  &=\int_{\overline M_{0,n+1}}\prod_{i=1}^{n}
       \pi_{n+1}^{*}\psi_{i,n}^{d_i}\\
  &=\int_{\overline M_{0,n+1}}\prod_{d_i>0}
       \bigl(\psi_{i,n+1}^{d_i}+(-e_{i,n+1})^{d_i}\bigr)\\
  &=\int_{\overline M_{0,n+1}}\prod_{i=1}^{n}\psi_{i,n+1}^{d_i}
    +\sum_{d_i>0}\int_{\overline M_{0,n+1}}
       (-e_{i,n+1})^{d_i}\prod_{j\ne i}
       \pi_{n+1}^{*}\psi_{j,n}^{d_j}\\
  &=\int_{\overline M_{0,n+1}}\prod_{i=1}^{n}\psi_{i,n+1}^{d_i}
    -\sum_{d_i>0}\int_{Z_{i,n+1}}
       (-e_{i,n+1})^{d_i-1}\prod_{j\ne i}
       \pi_{n+1}^{*}\psi_{j,n}^{d_j}\\
  &=\int_{\overline M_{0,n+1}}\prod_{i=1}^{n}\psi_{i,n+1}^{d_i}
    -\sum_{d_i>0}\int_{\overline M_{0,n}}
       \psi_{i,n}^{d_i-1}\prod_{j\ne i}\psi_{j,n}^{d_j}.
\end{align*}
第四个等号利用 $e_{i,n+1}$ 是子流形 $Z_{i,n+1}$ 的 Thom 类；最后一个等号则通过全纯微分同胚 $\sigma_{i,n}:\overline M_{0,n}\to Z_{i,n+1}$ 拉回被积式，由积分的度数定理得到。这就证明了~\eqref{eq:D-6-9}。

等式~\eqref{eq:D-6-9} 可写为
\begin{equation}\label{eq:D-6-10}
  \int_{\overline M_{0,n+1}}\prod_{i=1}^{n}\frac{1}{1-t_i\psi_i}
  =(t_1+\cdots+t_n)
    \int_{\overline M_{0,n}}\prod_{i=1}^{n}\frac{1}{1-t_i\psi_i}.
\end{equation}
这是\emph{\BookRecordIndex{T0345}弦方程}的一个特例。这里又把 $\psi_{i,n}$ 或 $\psi_{i,n+1}$ 简记为 $\psi_i$，第二下标可由上下文确定。

\emph{第三步。} 证明~\eqref{eq:D-6-6}。对 $n$ 归纳：当 $n=3$ 时，$\overline M_{0,3}$ 是一个点，结论显然成立。设对某个 $n\ge3$，\eqref{eq:D-6-6} 已成立，于是~\eqref{eq:D-6-7} 对该 $n$ 成立。由\BookRecordIndex{T0345}弦方程~\eqref{eq:D-6-10} 可得
\[
  \int_{\overline M_{0,n+1}}\prod_{i=1}^{n}\frac{1}{1-t_i\psi_i}
  =(t_1+\cdots+t_n)^{n-2}.
\]
因此，若非负整数 $d_1,\ldots,d_n$ 满足 $d_1+\cdots+d_n=n-2$，则
\begin{equation}\label{eq:D-6-11}
  \int_{\overline M_{0,n+1}}
    \psi_1^{d_1}\smile\psi_2^{d_2}\smile\cdots\smile\psi_n^{d_n}
  =\frac{(n-2)!}{d_1!\cdots d_n!}.
\end{equation}
现在令非负整数 $d_1,\ldots,d_{n+1}$ 满足 $d_1+\cdots+d_{n+1}=n-2$。其中必有一个 $d_i=0$。交换指标后由~\eqref{eq:D-6-11} 得
\begin{equation}\label{eq:D-6-12}
  \int_{\overline M_{0,n+1}}
    \psi_1^{d_1}\smile\psi_2^{d_2}\smile\cdots\smile\psi_{n+1}^{d_{n+1}}
  =\frac{(n-2)!}{d_1!\cdots d_{n+1}!}.
\end{equation}
这就完成了第三步及定理的证明。
\end{proof}
